% !TEX program = pdflatex
% First structured draft for the RaFT-UAV / RaFT-UAV++ paper.
% This file is intentionally bullet-point-ish: it is meant as a scaffold for a
% fast conference draft, not as final prose.

\documentclass[10pt,conference]{IEEEtran}

\usepackage{amsmath}
\usepackage{booktabs}
\usepackage{graphicx}
\usepackage{siunitx}
\usepackage{xcolor}
\usepackage{xspace}
\usepackage{enumitem}
\usepackage[hidelinks]{hyperref}

\sisetup{detect-all,per-mode=symbol}
\newcommand{\raftpp}{RaFT-UAV++\xspace}
\newcommand{\aerpaw}{AERPAW Dataset-28\xspace}
\newcommand{\todo}[1]{\textcolor{red}{[TODO: #1]}}
\newcommand{\figplaceholder}[2]{%
  \fbox{\begin{minipage}[c][#2][c]{0.92\linewidth}\centering\small #1\end{minipage}}%
}

\title{\raftpp: Artifact-Aware RF/Radar Fusion for UAV Tracking on AERPAW Dataset-28}

\author{%
  \IEEEauthorblockN{Author One, Author Two, Author Three}
  \IEEEauthorblockA{Affiliation(s) redacted for draft\\
  Email: \texttt{TODO@example.edu}}
}

\begin{document}
\maketitle

\begin{abstract}
\begin{itemize}[leftmargin=*,nosep]
  \item Problem: published RF/radar UAV tracking baselines on \aerpaw are hard to reproduce from raw streams because dataset artifacts, MATLAB post-processing, radar target selection, RF filtering, and evaluation windows are under-specified.
  \item We separate two questions that are often conflated: (i) Table-II-style artifact parity and (ii) raw-stream UAV tracking from RF CSV + Fortem radar JSON.
  \item We reproduce the paper-style Opt2 MATLAB artifact statistics and show why these artifact vectors are not equivalent to direct raw-stream replay.
  \item We introduce \raftpp, a truth-free raw-stream tracker combining global radar tracklet association, bounded unsupported-RF soft anchoring, same-track post-gap radar reacquisition, terminal short-segment guarding, and first-selected-radar bootstrap.
  \item On the repository raw-stream benchmark over Opt1--Opt3, \raftpp improves pooled 3D mean/p95/max from \SI{85.79}/\SI{398.03}/\SI{748.55}{m} to \SI{59.37}/\SI{106.21}/\SI{573.22}{m}.
\end{itemize}
\end{abstract}

\begin{IEEEkeywords}
UAV tracking, RF sensing, radar tracking, sensor fusion, data association, benchmark reproducibility, AERPAW, fixed-lag smoothing.
\end{IEEEkeywords}

\section{Introduction}
\label{sec:introduction}

\subsection{Problem statement}
\begin{itemize}[leftmargin=*]
  \item UAV tracking from passive/cellular RF localization and Fortem radar is attractive because each modality compensates for the other's failure modes.
  \item In practice, published results depend heavily on:\vspace{-0.25em}
  \begin{itemize}[leftmargin=*]
    \item which RF file/variant is used;
    \item how timestamps are corrected;
    \item whether radar positions come from Fortem LLA, raw polar fields, MATLAB-rotated tracks, or precomputed error vectors;
    \item how Fortem track IDs are selected and stitched;
    \item whether metrics are raw measurement errors, posterior-state errors, or artifact/proxy rows.
  \end{itemize}
  \item Our initial attempts with direct raw JSON/CSV replay could not match the reference paper, motivating a forensic split between paper-artifact parity and raw-stream tracking.
\end{itemize}

\subsection{Contributions}
\begin{enumerate}[leftmargin=*]
  \item \textbf{Artifact-level reproduction.} We identify and reproduce the Opt2 MATLAB artifact path used for paper-style RF/radar standalone statistics.
  \item \textbf{Raw-stream benchmark protocol.} We define a non-oracle raw-stream comparison over Opt1--Opt3 using RF CSV, Fortem radar JSON, and truth only for evaluation.
  \item \textbf{\raftpp method.} We propose global tracklet association with bounded unsupported-RF updates and radar-path bootstrap.
  \item \textbf{Failure analysis.} We show that Opt3's previous catastrophic failure is primarily an initialization/support problem, not a missing graph path or global RF clock offset.
  \item \textbf{Reproducible result.} \raftpp improves pooled raw-stream 3D mean/p95/max from \SI{85.79}/\SI{398.03}/\SI{748.55}{m} to \SI{59.37}/\SI{106.21}/\SI{573.22}{m}.
\end{enumerate}

\begin{figure}[t]
  \centering
  \figplaceholder{Overview placeholder: raw RF CSV + Fortem radar JSON $\rightarrow$ artifact audit and raw-stream \raftpp pipeline. Include two clearly separated lanes: ``paper-artifact parity'' and ``raw-stream method''.}{3.5cm}
  \caption{Proposed paper structure and evaluation split. Artifact parity explains the reference-paper statistics; \raftpp is evaluated separately on raw RF/radar streams.}
  \label{fig:overview}
\end{figure}

\section{Dataset and Reproducibility Pitfalls}
\label{sec:dataset}

\subsection{AERPAW Dataset-28 / AADM2025Dryad}
\begin{itemize}[leftmargin=*]
  \item Flights used in raw-stream benchmark: Opt1, Opt2, Opt3.
  \item Inputs for raw-stream tracker:\vspace{-0.25em}
  \begin{itemize}[leftmargin=*]
    \item Keysight RF localization CSV files;
    \item Fortem radar JSON track candidates;
    \item vehicleOut ground truth files used \emph{only} for post-run evaluation.
  \end{itemize}
  \item Table-II artifact path uses a distinct Opt2 MATLAB post-processing route; this is \emph{not} the same as raw JSON/CSV tracking.
\end{itemize}

\subsection{Artifact parity versus raw-stream tracking}
\begin{itemize}[leftmargin=*]
  \item Artifact parity reproduces paper-style RF/radar error vectors and proxy fusion statistics.
  \item Raw-stream tracking must operate only on raw RF/radar streams and truth-free association/scoring.
  \item We report these separately to avoid claiming that artifact/proxy rows are direct filter reproduction.
\end{itemize}

\begin{figure}[t]
  \centering
  \figplaceholder{Placeholder for a data-flow diagram: Opt2 MATLAB artifact path (Opt2.csv, -4 h RF correction, drop first 25 RF samples, radarerrorstandalone.mat) versus raw-stream Opt1--Opt3 path.}{3.25cm}
  \caption{Data-path ambiguity. The MATLAB artifact path reproduces Table-II-style statistics, while the raw-stream path uses RF CSV and Fortem JSON directly.}
  \label{fig:artifact-vs-raw}
\end{figure}

\section{Paper-Artifact Forensics}
\label{sec:artifact_forensics}

\subsection{Opt2 MATLAB artifact path}
\begin{itemize}[leftmargin=*]
  \item RF artifact source:\vspace{-0.25em}
  \begin{itemize}[leftmargin=*]
    \item \texttt{RF Sensor and Radar/Post\_Processing\_Scripts/Opt2.csv}
    \item subtract \SI{4}{h} from RF timestamps;
    \item nearest-timestamp truth matching;
    \item drop first 25 RF samples;
    \item keep finite nonzero fixes.
  \end{itemize}
  \item Radar artifact source:\vspace{-0.25em}
  \begin{itemize}[leftmargin=*]
    \item \texttt{radarerrorstandalone.mat}
    \item synthetic MATLAB radar timing $t_0=\SI{0.25}{s}$, $T_s=\SI{0.25}{s}$.
  \end{itemize}
  \item Actual CV/fixed-lag replay over \texttt{radar\_outputOpt2.mat} is diagnostic and differs from the artifact vector.
\end{itemize}

\begin{table}[t]
\centering
\caption{Artifact-level parity summary. Fill exact values from \texttt{paper\_table\_ii\_artifact\_opt2\_conclusion\_table.csv}.}
\label{tab:artifact_parity}
\begin{tabular}{lrrrr}
\toprule
Row & Count & Mean [m] & Std/P95 [m] & Max [m] \\
\midrule
RF raw artifact & 206 & 471.75 & 885.24 & 4831.63 \\
RF after NIS selected & 125 & \todo{25--28} & \todo{} & \todo{} \\
Radar raw artifact & 3106 & 26.19 & \todo{} & 195.57 \\
Radar after NIS artifact & 2403 & 20.95 & \todo{} & 97.20 \\
KF all-steps proxy & 2655 & 22.09 & \todo{} & \todo{} \\
Actual CV replay diagnostic & \todo{} & 25.20 & 82.12 & 102.20 \\
\bottomrule
\end{tabular}
\end{table}

\subsection{Key forensic finding}
\begin{itemize}[leftmargin=*]
  \item The artifact error vectors are reproducible, but the row-level radar coordinate stream used for actual replay is not identical to the artifact vector.
  \item This explains why artifact parity is achievable while direct raw replay cannot be treated as exact Table-II fusion reproduction.
  \item Paper claim should be: \emph{artifact-level statistics are reproduced; raw-stream tracking is evaluated separately.}
\end{itemize}

\section{RaFT-UAV++ Method}
\label{sec:method}

\subsection{Overview}
\begin{itemize}[leftmargin=*]
  \item State: constant-velocity 3D state $x_t=[e,n,u,\dot e,\dot n,\dot u]^\top$.
  \item Radar provides 3D ENU candidate track points, but altitude is treated weakly or with large covariance.
  \item RF provides 2D position fixes with highly variable reliability; unsupported RF is not fully trusted and not fully rejected.
  \item Fixed-lag smoothing with lag \SI{20}{s} is used in the final raw-stream method.
\end{itemize}

\begin{figure}[t]
  \centering
  \figplaceholder{Algorithm placeholder: global segment DAG $\rightarrow$ top-k path generation $\rightarrow$ replay ranking $\rightarrow$ first-selected-radar bootstrap $\rightarrow$ bounded RF soft-anchor $\rightarrow$ fixed-lag smoothing.}{4.0cm}
  \caption{\raftpp raw-stream method. The graph is selected without truth; truth is used only for evaluation.}
  \label{fig:method_pipeline}
\end{figure}

\subsection{Global radar tracklet graph}
\begin{itemize}[leftmargin=*]
  \item Split Fortem radar candidates into continuous same-track-ID microsegments.
  \item Segment costs include:\vspace{-0.25em}
  \begin{itemize}[leftmargin=*]
    \item motion plausibility;
    \item internal speed/range discontinuities;
    \item RF contradiction;
    \item catProb/confidence when available;
    \item path-switch penalties;
    \item terminal short-segment penalties.
  \end{itemize}
  \item A diverse top-$K$ path set is generated and replay-ranked by truth-free innovation/residual summaries.
\end{itemize}

\subsection{First-selected-radar bootstrap}
\begin{itemize}[leftmargin=*]
  \item Problem: early RF/coast records can initialize the filter far from the radar-supported path, especially in Opt3.
  \item Fix: after graph association, initialize replay from the first graph-selected radar event.
  \item Pre-bootstrap events are dropped/recorded, not assimilated silently.
  \item Optional variants considered but not final: first selected radar with local-linear velocity, first radar after 30 s.
\end{itemize}

\begin{figure}[t]
  \centering
  \figplaceholder{Opt3 bootstrap placeholder: plot error over time before/after first-selected-radar bootstrap, highlighting early coast/unsupported RF blow-up and its removal.}{3.25cm}
  \caption{First-selected-radar bootstrap removes the Opt3 early-state failure that dominated the previous raw-stream method.}
  \label{fig:opt3_bootstrap}
\end{figure}

\subsection{Unsupported-RF soft-anchor}
\begin{itemize}[leftmargin=*]
  \item Unsupported RF is mixed quality: some fixes are useful; others have kilometer-scale errors.
  \item Binary choices fail:\vspace{-0.25em}
  \begin{itemize}[leftmargin=*]
    \item fully trusting unsupported RF causes drift/jumps;
    \item rejecting/coasting through unsupported RF can also drift.
  \end{itemize}
  \item Soft-anchor policy:\vspace{-0.25em}
  \begin{itemize}[leftmargin=*]
    \item predict to RF time;
    \item compute 2D innovation;
    \item cap correction length at \SI{5}{m};
    \item apply conservative covariance;
    \item leave vertical state policy as \texttt{keep} in final method.
  \end{itemize}
\end{itemize}

\subsection{Post-gap radar bridge and terminal guard}
\begin{itemize}[leftmargin=*]
  \item Same-track post-gap radar bridge allows short radar reacquisitions after unsupported-RF blocks.
  \item Bridge updates are allowed only for eligible same-track post-gap rows.
  \item Terminal short-segment guard penalizes unsupported late terminal segments, preventing Opt1-style track-2351 tails.
\end{itemize}

\section{Experimental Protocol}
\label{sec:experiments}

\subsection{Raw-stream benchmark}
\begin{itemize}[leftmargin=*]
  \item Flights: Opt1, Opt2, Opt3.
  \item Metrics:\vspace{-0.25em}
  \begin{itemize}[leftmargin=*]
    \item truth-grid 3D mean, RMSE, p95, max;
    \item paper-style nearest-truth sample-time metrics;
    \item selected radar rows;
    \item accepted RF rows;
    \item track switches;
    \item unsupported RF count and max.
  \end{itemize}
  \item Paper-artifact rows are excluded from raw-stream comparison tables.
\end{itemize}

\subsection{Baselines}
\begin{itemize}[leftmargin=*]
  \item CV catProb fixed-lag baseline.
  \item Prediction-NIS fixed-lag baseline.
  \item Range-covariance tracklet-Viterbi fixed-lag baseline.
  \item Previous best non-oracle preset.
  \item \raftpp final method.
\end{itemize}

\begin{table*}[t]
\centering
\caption{Raw-stream pooled comparison over Opt1--Opt3. Artifact/proxy rows are excluded.}
\label{tab:raw_pooled}
\begin{tabular}{lrrrr}
\toprule
Method & Mean 3D [m] & RMSE 3D [m] & P95 3D [m] & Max 3D [m] \\
\midrule
CV catProb fixed-lag & 109.88 & 226.32 & 410.97 & 2098.57 \\
Prediction-NIS fixed-lag & 169.94 & 493.09 & 471.11 & 5108.39 \\
Range-covariance tracklet-Viterbi & 93.06 & 177.67 & 391.51 & 1744.58 \\
Best previous non-oracle preset & 85.79 & 130.98 & 398.03 & 748.55 \\
\textbf{\raftpp} & \textbf{59.37} & \todo{fill} & \textbf{106.21} & \textbf{573.22} \\
\bottomrule
\end{tabular}
\end{table*}

\begin{table}[t]
\centering
\caption{Per-flight \raftpp result.}
\label{tab:raftpp_by_flight}
\begin{tabular}{lrrrr}
\toprule
Flight & Mean 3D [m] & P50 [m] & P95 [m] & Max [m] \\
\midrule
Opt1 & 57.92 & 51.79 & 106.93 & 573.22 \\
Opt2 & 57.16 & 48.17 & 126.60 & 290.47 \\
Opt3 & 62.44 & 56.51 & 85.71 & 334.75 \\
\bottomrule
\end{tabular}
\end{table}

\section{Ablation Plan}
\label{sec:ablation}

\begin{itemize}[leftmargin=*]
  \item Required before final paper submission: generate the table below from clean committed runs.
  \item Important rows:\vspace{-0.25em}
  \begin{itemize}[leftmargin=*]
    \item previous best non-oracle preset;
    \item global tracklet without soft-anchor;
    \item soft-anchor cap5 without post-gap/terminal logic;
    \item post-gap bridge without terminal guard;
    \item terminal guard without first-radar bootstrap;
    \item full \raftpp.
  \end{itemize}
\end{itemize}

\begin{table*}[t]
\centering
\caption{Ablation placeholder. Fill from clean all-flight reruns.}
\label{tab:ablation}
\begin{tabular}{lrrrrrr}
\toprule
Method variant & Mean 3D & P95 3D & Max 3D & Unsupported RF max & Track switches & Notes \\
\midrule
Previous best non-oracle & 85.79 & 398.03 & 748.55 & \todo{} & \todo{} & baseline \\
Global tracklet only & \todo{} & \todo{} & \todo{} & \todo{} & \todo{} & path graph \\
+ soft-anchor cap5 & \todo{} & \todo{} & \todo{} & \todo{} & \todo{} & bounded RF \\
+ post-gap bridge & \todo{} & \todo{} & \todo{} & \todo{} & \todo{} & reacquisition \\
+ terminal guard & \todo{} & \todo{} & \todo{} & \todo{} & \todo{} & tail control \\
+ first-radar bootstrap (\raftpp) & 59.37 & 106.21 & 573.22 & \todo{} & \todo{} & final \\
\bottomrule
\end{tabular}
\end{table*}

\section{Discussion}
\label{sec:discussion}

\subsection{What was learned from artifact parity}
\begin{itemize}[leftmargin=*]
  \item Table-II artifact vectors can be reproduced, but this is not the same as raw-stream Kalman replay.
  \item Hidden post-processing, selected radar targets, timestamp correction, and MATLAB artifacts materially affect reported results.
\end{itemize}

\subsection{What was learned from raw-stream tracking}
\begin{itemize}[leftmargin=*]
  \item Unsupported RF should neither be fully trusted nor fully rejected; bounded soft anchoring is effective.
  \item Radar association must be global/path-level rather than frame-local.
  \item Bootstrap from the graph-selected radar path fixes early-state failures, especially in Opt3.
  \item Terminal short radar segments can harm p95/mean when unsupported by RF or continuity.
\end{itemize}

\subsection{Limitations}
\begin{itemize}[leftmargin=*]
  \item Method was developed on Dataset-28; if no LOFO scalar validation is included, avoid claiming held-out generalization.
  \item Fixed-lag smoothing introduces latency; report offline/fixed-lag status clearly.
  \item Artifact parity and raw-stream method are separate claims.
  \item Future work: recover the hidden radar coordinate stream behind \texttt{radarerrorstandalone.mat}; evaluate additional AERPAW flights; test online-only variants.
\end{itemize}

\section{Conclusion}
\label{sec:conclusion}

\begin{itemize}[leftmargin=*]
  \item We provide a reproducible separation between paper-artifact parity and raw-stream RF/radar UAV tracking.
  \item We introduce \raftpp, a global-tracklet fixed-lag tracker with bounded RF anchoring and radar-path bootstrap.
  \item On Opt1--Opt3, \raftpp improves pooled raw-stream 3D mean/p95/max from \SI{85.79}/\SI{398.03}/\SI{748.55}{m} to \SI{59.37}/\SI{106.21}/\SI{573.22}{m}.
\end{itemize}

\section*{Checklist before submission}
\begin{itemize}[leftmargin=*]
  \item Regenerate all tables from a clean commit.
  \item Fill ablation table.
  \item Add method pseudocode as Algorithm 1 if page budget allows.
  \item Add trajectory/error-over-time figures for Opt1--Opt3.
  \item Add artifact parity table and short data-flow figure.
  \item Add non-oracle guardrail tests for first-selected-radar bootstrap and path scoring.
  \item Decide exact target venue format: IEEE 6-page conference or longer journal version.
\end{itemize}

\begin{thebibliography}{9}

\bibitem{aerpaw_dataset28}
AERPAW / Dryad Dataset-28 authors,
``Multi-modal RF Sensor and Radar Dataset for UAV Tracking,''
Dryad DOI: \texttt{10.5061/dryad.7d7wm3898},
\todo{complete citation}.

\bibitem{reference_paper}
\todo{Authors},
``Fusion of Cellular ISAC and Passive RF Sensing for UAV Detection and Tracking,''
\todo{venue, year}.

\bibitem{kalman}
R.~E. Kalman,
``A New Approach to Linear Filtering and Prediction Problems,''
\todo{complete citation}.

\bibitem{tracklet_data_association}
\todo{Add representative tracklet / multi-target tracking / data association references.}

\bibitem{robust_sensor_fusion}
\todo{Add robust sensor fusion / NIS / Student-t / Huber references if used in final comparison.}

\end{thebibliography}

\end{document}
